% vim: ai sw=2

\documentclass[xcolor=dvipsnames,handout]{beamer}

\let\Oldemph\emph
\renewcommand{\emph}[1]{\textcolor{blue}{\Oldemph{#1}}}
\newcommand{\heading}[1]{\textbf{\textcolor{BurntOrange}{#1}}}

% Copyright 2004 by Till Tantau <tantau@users.sourceforge.net>.
%
% In principle, this file can be redistributed and/or modified under
% the terms of the GNU Public License, version 2.
%
% However, this file is supposed to be a template to be modified
% for your own needs. For this reason, if you use this file as a
% template and not specifically distribute it as part of a another
% package/program, I grant the extra permission to freely copy and
% modify this file as you see fit and even to delete this copyright
% notice. 


\mode<presentation>
{
%  \usetheme{Warsaw}
\usetheme{Madrid}

  \setbeamercovered{transparent}
}

\usepackage[utf8]{inputenc}
\usepackage{times}
\usepackage[T1]{fontenc}

\usepackage{graphicx}
\usepackage[final]{listings}
\usepackage{multirow}

\lstset{tabsize=2,numbers=none,showstringspaces=false,breaklines=true,
	breakatwhitespace=true,escapeinside={(*@}{@*)},
	basicstyle=\scriptsize,keywordstyle=\bfseries\color{OliveGreen},
	commentstyle=\color{blue}}

\title[PB173/01] % (optional, use only with long paper titles)
{PB173 -- Ovladače jádra -- Linux}

\subtitle{I.}

\author[J. Slabý]{Jiří~Slabý}

\institute[ITI] % (optional, but mostly needed)
{ITI, Fakulta Informatiky}

\date{21. 9. 2010}

\subject{}
% This is only inserted into the PDF information catalog. Can be left
% out. 



% If you have a file called "university-logo-filename.xxx", where xxx
% is a graphic format that can be processed by latex or pdflatex,
% resp., then you can add a logo as follows:

% \pgfdeclareimage[height=0.5cm]{university-logo}{university-logo-filename}
% \logo{\pgfuseimage{university-logo}}



% Delete this, if you do not want the table of contents to pop up at
% the beginning of each subsection:
% \AtBeginSubsection[]
% {
%   \begin{frame}<beamer>
%     \frametitle{Outline}
%     \tableofcontents[currentsection,currentsubsection]
%   \end{frame}
% }


% If you wish to uncover everything in a step-wise fashion, uncomment
% the following command: 

%\beamerdefaultoverlayspecification{<+->}

\begin{document}

\begin{frame}
  \titlepage
\end{frame}

%%%%%%%%%%%%%%%%%%%%%%%

\begin{frame}
  \frametitle{Úvodní informace}
  \begin{itemize}
  \item Semestr = 13 týdnů (28. 9. svátek)
  \item Cvičící
  \begin{itemize}
  \item Vývoj jádra od r. 2005
  \item Student FI
  \end{itemize}
  \item Cíle cvičení
  \begin{itemize}
  \item Nastínit trochu jiný model programování
  \item Prohloubit znalosti vnitřností OS a HW
  \end{itemize}
  \item Ukončení: k
  \begin{itemize}
  \item Splnění úkolů (možno ve skupinkách)
  \end{itemize}
  \end{itemize}

  Vše potřebné na:
  {\small\url{https://minotaur.fi.muni.cz:8443/~xsvenda/docuwiki/doku.php?id=public:pb173:pb173_2010_systemlinux}}

  \bigskip
  \pause

  \begin{center}
  \heading{Vstupní znalosti?}
  \end{center}
\end{frame}

\begin{frame}
  \frametitle{S čím budeme pracovat?}
  \begin{itemize}

  \item satyr01-satyr10
  \begin{itemize}
  \item CentOS 5.5
  \end{itemize}

  \item GIT
  \begin{itemize}
  \item Úvod do GITu zde
  \item Podrobněji: \url{http://book.git-scm.com/}
  \end{itemize}

  \item Zdroje jádra
  \begin{itemize}
  \item Použijeme předinstalované z RPM (\texttt{/usr/src/kernels/})
  \item \url{http://git.kernel.org/}
  \end{itemize}

  \item \textsc{Qemu}
  \begin{itemize}
  \item Ze začátku
  \end{itemize}

  \item \textsc{Combo6X} karty
  \begin{itemize}
  \item Později; k PCI, I/O, na přerušení, DMA
  \item \url{http://www.liberouter.org/}
  \end{itemize}

  \end{itemize}
\end{frame}

\begin{frame}
  \frametitle{Úkol první}
  \heading{Práce s GITem}
  \begin{enumerate}
  \item \texttt{git clone git://decibel.fi.muni.cz/\textasciitilde xslaby/pb173}

  \pause

  \item Změňte obsah souboru \texttt{sandbox/hello}
  \item Zkontrolujte změny (\texttt{git diff --color})
  \item \texttt{git commit}

  \pause

  \item Smažte \texttt{sandbox/hello} (\texttt{git rm})
  \item \texttt{git commit}

  \pause

  \item Vygenerujte 2 patche (\texttt{git format-patch -2})
  \end{enumerate}
\end{frame}

\begin{frame}
  \frametitle{Záplatování}

  \heading{Co s patchem dále?}
  \begin{itemize}
  \item Poslat odpovídajícímu správci (\texttt{scripts/get\_maintainer.pl})
  \item Poslat na ML a vystavit celý GIT strom, aby někdo provedl ,,merge''
  \end{itemize}

  \heading{Vývojový proces jádra}
  \begin{itemize}
  \item L. Torvalds má právo veta
  \item Záplatováním se opravuje i přidává nová funkcionalita
  \item Ovladače mohou být zařazeny přes \texttt{drivers/staging}
  \begin{itemize}
  \item Začištění, smazání nepotřebného a přesun ze staging
  \item Správce Greg KH
  \end{itemize}
  \end{itemize}

  \bigskip
  \pause

  \begin{center}
  \heading{Stejným způsobem odevzdávání úkolů}
  \end{center}
\end{frame}

\begin{frame}
  \frametitle{Jádro vs. programy}

  \heading{Hlavní rozdíly}
  \begin{itemize}
  \item Žádné \emph{libc} (printf, strlen, malloc, \ldots), ani ostatní
    (pthread)
  \item \emph{Paměťový prostor}
  \item \emph{Počáteční funkce} (main)
  \item Pád systému $\Rightarrow$ pád všeho
  \end{itemize}
\end{frame}

\begin{frame}
  \frametitle{Kód jádra}

  \begin{itemize}
  \item GNU C (x86 už jen \texttt{gcc} 4.x)
  \item \emph{CodingStyle}
  \begin{itemize}
  \item Kontrola: \texttt{scripts/checkpatch.pl} (ne 100\%)
  \end{itemize}
  \item Monolit (ale moduly)
  \item \texttt{module\_init}
  \end{itemize}
\end{frame}

\begin{frame}
  \frametitle{Úkol druhý}

  \begin{itemize}
  \item ss
  \end{itemize}
\end{frame}

\begin{frame}[fragile]
  \frametitle{Configuration}
  \heading{Automaton Checker configuration example}
  \begin{lstlisting}[language=C]
int a;
  \end{lstlisting}
\end{frame}

\end{document}

